\documentclass[11pt]{article}
\usepackage[T1]{fontenc}
\usepackage{lmodern,amsmath,amssymb,amsthm,booktabs}
\usepackage[margin=1in]{geometry}
\usepackage[hidelinks]{hyperref}
\newtheorem{proposition}{Proposition}
\newtheorem{corollary}[proposition]{Corollary}
\title{Direct axis routing removes a cubic exponent loss}
\author{Douglas Colkitt\thanks{Prepared with assistance from OpenAI Codex. This is not independent review or endorsement by OpenAI.}}
\date{October 7, 2026}
\begin{document}
\maketitle
\begin{abstract}
We audit a scheduling improvement to OpenAI's integer multiplication manuscript.
Its existing arbitrary-gap interchange primitive permits the CRT axis reversal
and resampling axis exposure to use $O(d)$ field interchanges rather than
$O(d^2)$ adjacent interchanges. The resulting layout cost removes the constraint
that the dimension exponent must be as small as the network saving. Conditional
on the upstream algorithmic interfaces, exact parameters give
$T(n)=O(n(\log n)^{1-\kappa})$ with $\kappa=2^{-78}$ using the previously
audited $h=46$ network. Keeping the original network and recurrence exponents
instead gives $\kappa=2^{-107}$. The earlier cubic ceiling concerns the original
cost accounting and does not apply to the improved schedules.
\end{abstract}

\section{Scope and primitive}
The source is OpenAI's \emph{Integer multiplication below $n\log n$}, pinned at
commit
\begin{center}
\texttt{adc7f1241b42e322a6451854ab7e4b4c146bf78a}.
\end{center}
We use its fixed
finite-alphabet Turing-machine model with a fixed finite number of
one-dimensional tapes. The source algorithmic interfaces, including their
uniform time bounds, remain assumptions of this conditional improvement.

Let $V$ be the stored bit volume of a complete padded tensor with axes of
lengths $t_i\in\{2^{\ell-1},2^\ell\}$ and complete coefficient records.
Lemma \texttt{lem:chunk-swap} accepts the address shape
\[
 [P]\times[2^u]\times[G]\times[2^u]\times[B]
\]
and interchanges the two width-$u$ fields in $O(Vu^\tau)$ steps. Prefix,
gap, and suffix lengths are arbitrary positive integers; $B$ may include the
entire coefficient word. The paragraph following the lemma already extends
this to widths differing by one. Writing $e$ for the extra bit gives
\[
 (e,x,G,y)\longmapsto(e,y,G,x)\longmapsto(y,G,e,x),
\]
\[
 (x,G,e,y)\longmapsto(e,x,G,y)\longmapsto(e,y,G,x),
\]
where $x,y$ have equal width. Each case costs one interchange and one
$O(V)$ single-bit move. Thus a nonadjacent full-field swap costs
$O(V(1+\ell^\tau))$. Empty equal-width pieces require no swap.
These are permutations of complete records; they preserve every record's
internal bit order. No new network or interchange primitive is introduced.

\section{The two schedules}
\begin{proposition}[Linear count of field interchanges]
Conditional on the stated stream and swap primitives, CRT axis reversal and
exposure/restoration of every resampling axis cost
$O(Vd(1+\ell^\tau))$, including descriptors and cleanup.
\end{proposition}
\begin{proof}
For the CRT map, the padded input order is $(a_d,\ldots,a_1)$.
Swap positions $j$ and $d+1-j$ for $1\le j\le\lfloor d/2\rfloor$.
The disjoint swaps produce exactly $(a_1,\ldots,a_d)$. In the existing
triangular formula
\[
 b_i=a_i+\mu_i\sum_{j<i}a_jP_j\pmod{s_i},\qquad
 P_i=\prod_{j<i}s_j,\quad \mu_i=P_i^{-1}\pmod{s_i},
\]
targets are still processed in order $d,d-1,\ldots,1$. All controls precede
the target and still contain the required original digits. The inverse
recovers digits in order $1,\ldots,d$ and then reverses the field permutation.
Padding insertion precedes the forward reversal; deletion follows the inverse.
The original enumerations are therefore unchanged.

For resampling, swap the next named axis with the innermost field, process
its lines, and undo that swap. An already innermost axis needs no swap.
This uses at most $2(d-1)$ field interchanges per tensor map. Axis names,
widths, and validity bounds move with their fields. A line fixes all other
named coordinates and traverses the target in increasing numerical order.
Their temporary physical order does not affect its validity predicate.

Keep the original axis-processing order and restore the original physical
field order after every axis. The resulting operator is exactly the original
coordinatewise map, with the same line arithmetic and truncations. Expansion
writes all $t_i$ outputs; compression writes $s_i$ outputs and pads with zeros.
The full box stays allocated. Invalid entries remain zero at the completed
axis-map boundaries. An intermediate bit move or the internal arithmetic of a
swap need not preserve an intermediate physical validity mask: only the
completed permutation is used by the next numerical operation. Exact
permutations preserve norms, so the contraction and error bounds are unchanged.

Each call collapses all intervening axes into its single spectator field.
Since $t_i\ge2$, $d\le\log_2 T$, where $T=\prod_i t_i$. The named-axis
list, record-width description, and all collapsed lengths have size polynomial
in $\log(2V)$. A call descriptor can be computed by ordinary integer arithmetic
in $\operatorname{poly}(\log(2V))=O(V)$ time. Thus list scans and updates are
charged at every interchange. The primitives run sequentially, clean their
workspaces, and reuse a fixed tape collection. All growing lists are stored
data. Multiplying the per-interchange bound by the linear count proves the claim.
\end{proof}

The $d$ controlled CRT rotations cost $O(dV)$ as before. Numerical resampling
costs and per-line setup do not change. In particular, no assertion that
approximate maps commute is needed: the axis-processing order is preserved.

\section{Changed accounting and explicit witnesses}
The source has $p=\Theta(\log n)$, $V=\Theta(n)$,
$d=\Theta(p^\epsilon)$ and $\ell=\Theta(p^{1-\epsilon})$. The old layout
term and its exponent were
\[
 d^2(1+\ell^\tau),\qquad \tau+\epsilon(2-\tau).
\]
The proposition replaces these by
\[
 d(1+\ell^\tau),\qquad \tau+\epsilon(1-\tau).
\]
Consequently the fourth margin becomes
\[
 g_4=(1-\tau)(1-\epsilon).
\]
The other six margins retain their original formulas:
\begin{align*}
 g_1&=1-\epsilon(1+c),&g_2&=\epsilon c(1-\tau),\\
 g_3&=\epsilon(1-\lambda'),&g_5&=\tfrac14-\delta-\tfrac32\epsilon,\\
 g_6&=1-\delta-\epsilon,&g_7&=\epsilon.
\end{align*}
No analytic $d^2$ factor in the precision or prime-interval bounds is removed.

\begin{proposition}[Parameter substitution]
For either $a=2^{-50}$ with the original network, or
$a=9/500000000000$ with the previously audited $h=46$ network, set
\begin{gather*}
 \tau=\sigma=1-a,\quad \beta=\tfrac12,\quad
 \delta=\tfrac1{16},\quad C_1=20,\\
 \epsilon=\tfrac1{40},\quad c=a/2,\quad
 \lambda=1-3a^2/4,\quad \lambda'=1-a^2/2.
\end{gather*}
The modified assembly constraints and all recorded side conditions hold.
The minimum final margin is $G=a^2/80$.
\end{proposition}
\begin{proof}
The layer inequalities follow from
\[
 \tau(1+c/\beta)=(1-a)(1+a)=1-a^2<\lambda<\lambda'<1,
 \qquad 1-a/2<\lambda'.
\]
Also $\max\{\tau,\sigma\}<\lambda$. Both values of $a$ are less than
$1/16$, which is more than sufficient for these comparisons. The guard and
dimension conditions are $\epsilon C_1=1/2<1$ and $\epsilon<1/12$.
Direct substitution gives
\[
 g_2=g_3=\frac{a^2}{80},\quad
 g_4=\frac{39a}{40},\quad g_5=\frac3{20},\quad
 g_6=\frac{73}{80},\quad g_7=\frac1{40},\quad g_1>\frac12.
\]
The first two are the minima for either stated $a$.

The remaining dependencies can be checked explicitly:
\begin{center}
\begin{tabular}{ll}
\toprule
Quantity & Exponent or implication\\
\midrule
Guard $d^{20}$ & $p^{1/2}=o(p)$\\
Resampling $\alpha$ & $\Theta(p^{21/80})=o(p^{1/2})$\\
Resampling $\gamma$ & $O(p^{11/20})=o(p)$\\
Prime-interval ratio & $p^{1-2\epsilon}=p^{19/20}\to\infty$\\
Chunk $K$ & $\Theta(p^{a/80})$, hence $K/\log p\to\infty$\\
Axis width $\ell$ & $\Theta(p^{39/40})$, hence $K=o(\ell)$\\
Record length $r$ & $2^{\Theta(p^{39/40})}$ dominates every fixed polynomial\\
\bottomrule
\end{tabular}
\end{center}
Thus polynomially many exceptional repairs, descriptor preparations and
per-line setup retain their asymptotic bounds. Since $d\to\infty$, prime
search still takes $n^{o(1)}$ time. The explicit root calculation is the
fixed-power comparison $d^{40}\le b$. The original error bound still yields
correct rounding beyond one common cutoff because $\gamma=o(b)$ and $p=6b$.
\end{proof}

\begin{corollary}[Conditional multiplication bounds]
With the original network and recurrence exponents, the modified schedule
permits $\kappa=2^{-107}$. With the audited $h=46$ network and rational
recurrence saving, it permits $\kappa=2^{-78}$.
\end{corollary}
\begin{proof}
In the original-network case,
$G=2^{-100}/80>2^{-107}$, since $128>80$.
For $h=46$, exact arithmetic gives
\[
 G=\frac{81}{20000000000000000000000000}
   =4.05\cdot10^{-24}>2^{-78}.
\]
For the chosen $\kappa$ put $\rho=G-\kappa>0$. Every fixed power of
$\log p$ is eventually at most $p^\rho$, so the total normalized bound is
$O(p^{1-\kappa})$. Apply the unchanged final assembly and finite-cutoff argument.
\end{proof}

\section{Relationship to the earlier ceiling}
Under the original adjacent-layout accounting, positivity of
$g_4=a-\epsilon(1+a)$ forces $\epsilon<a$. The other layer inequalities
force $\min\{ac,1-\lambda'\}<a^2$, yielding $\kappa<a^3$.
The new $g_4=a(1-\epsilon)$ no longer forces $\epsilon<a$.
The explicit witnesses above are therefore outside the hypotheses of that
ceiling. They are rejected by the original numerical constraint checker and
accepted only when the nonadjacent layout model is explicitly selected.

This improvement retains the finite-network construction, butterfly recursion,
precision bound, and analytic reductions. It changes two routing schedules and
one row of cost accounting. It does not claim an optimal exponent for the new
accounting. Further tuning of $\epsilon$, the stopping exponent, or the guard
bound is a separate question.

\section{Artifacts and review boundary}
The \href{https://github.com/CrocSwap/integer-mult-bounds}{companion repository}
contains the note, certificates, and source patches.
This follow-up is supplied as a separate audit with three alternative patches:
\begin{itemize}
\item \texttt{nonadjacent-layout.patch}: the routing proof and cost changes only;
  two source files, original parameters retained.
\item \texttt{frozen-nonadjacent-107.patch}: four source files, retaining the
  original network, recurrence exponents, and $\beta=1/2$.
\item \texttt{h46-nonadjacent-78.patch}: seven source files, retaining the
  previously audited $h=46$ network and $\beta=1/2$.
\end{itemize}
All apply independently to the pinned upstream source. Run \texttt{make verify}
for exact certificates, the finite tests, source integrity, and patch application
checks. Run \texttt{make audit-note} to build this document.

Finite tests compare the schedules with independent tuple-coordinate oracles,
check padded CRT conversion and inversion, and compare routed resampling with
exact tensor-product sums. They do not establish the upstream swap routine's
Turing-machine time bound. The written argument supplies the new linear
schedule count and its composition; the unchanged upstream interfaces retain
their conditional status. Independent review of this extension and those
interfaces is still needed before an unconditional multiplication claim.
\end{document}
